\section[The propositional rules of LK]{The propositional rules of \(\LK\)}
\label{sec:lk-lj-02}

The initial and structural rules are
\[
\frac{}{A\Rightarrow A}\;\textsc{Id}
\qquad
\frac{}{\bottom,\Gamma\Rightarrow\Delta}\;\textsc{\(\bottom L\)},
\]
\[
\frac{\Gamma\Rightarrow\Delta}
     {A,\Gamma\Rightarrow\Delta}\;\textsc{WL}
\qquad
\frac{\Gamma\Rightarrow\Delta}
     {\Gamma\Rightarrow\Delta,A}\;\textsc{WR},
\]
\[
\frac{A,A,\Gamma\Rightarrow\Delta}
     {A,\Gamma\Rightarrow\Delta}\;\textsc{CL}
\qquad
\frac{\Gamma\Rightarrow\Delta,A,A}
     {\Gamma\Rightarrow\Delta,A}\;\textsc{CR},
\]
and
\[
\frac{\Gamma\Rightarrow\Delta,A
      \qquad
      A,\Pi\Rightarrow\Sigma}
     {\Gamma,\Pi\Rightarrow\Delta,\Sigma}\;\textsc{Cut}.
\]
The displayed occurrence of \(A\) in cut is the \emph{cut formula}.

The logical rules for conjunction and disjunction are
\[
\frac{A,B,\Gamma\Rightarrow\Delta}
     {A\land B,\Gamma\Rightarrow\Delta}\;\textsc{\(\land L\)}
\qquad
\frac{\Gamma\Rightarrow\Delta,A
      \qquad
      \Gamma\Rightarrow\Delta,B}
     {\Gamma\Rightarrow\Delta,A\land B}\;\textsc{\(\land R\)},
\]
\[
\frac{A,\Gamma\Rightarrow\Delta
      \qquad
      B,\Gamma\Rightarrow\Delta}
     {A\lor B,\Gamma\Rightarrow\Delta}\;\textsc{\(\lor L\)}
\qquad
\frac{\Gamma\Rightarrow\Delta,A,B}
     {\Gamma\Rightarrow\Delta,A\lor B}\;\textsc{\(\lor R\)}.
\]
The implication rules are
\[
\frac{\Gamma\Rightarrow\Delta,A
      \qquad
      B,\Pi\Rightarrow\Sigma}
     {A\to B,\Gamma,\Pi\Rightarrow\Delta,\Sigma}
     \;\textsc{\(\to L\)}
\qquad
\frac{A,\Gamma\Rightarrow\Delta,B}
     {\Gamma\Rightarrow\Delta,A\to B}
     \;\textsc{\(\to R\)}.
\]
Negation is the abbreviation \(\neg A\defeq A\to\bottom\).  If negation is
instead made primitive, its rules must be added and the cut-elimination proof
must include their cases.

\section{Quantifier rules}
\label{sec:lk-lj-03}

For first-order formulas the quantifier rules are
\[
\frac{A[t/x],\Gamma\Rightarrow\Delta}
     {\forall x\,A,\Gamma\Rightarrow\Delta}\;\textsc{\(\forall L\)}
\qquad
\frac{\Gamma\Rightarrow\Delta,A[y/x]}
     {\Gamma\Rightarrow\Delta,\forall x\,A}\;\textsc{\(\forall R\)},
\]
\[
\frac{A[y/x],\Gamma\Rightarrow\Delta}
     {\exists x\,A,\Gamma\Rightarrow\Delta}\;\textsc{\(\exists L\)}
\qquad
\frac{\Gamma\Rightarrow\Delta,A[t/x]}
     {\Gamma\Rightarrow\Delta,\exists x\,A}\;\textsc{\(\exists R\)}.
\]
Every displayed instance \(A[t/x]\) uses the total capture-avoiding
substitution fixed in the first-order syntax chapter; bound variables are
renamed when required, so there is no additional definedness side condition.
In \(\forall R\) and \(\exists L\), the eigenvariable \(y\) must not occur
free in the conclusion sequent.  This freshness requirement is a rule
condition, not a convention that may be repaired after the inference.

